\documentclass[10pt,a4paper]{amsart}
\usepackage[margin=19mm]{geometry}
\usepackage{amsmath,amssymb,mathtools}
\usepackage[scr=rsfs,cal=euler]{mathalfa}
\usepackage{xfrac}
\usepackage[unicode,hidelinks,bookmarks=false]{hyperref}
\newcommand{\ie}{i.e.\ }
\newcommand{\cf}{cf.\ }
\newcommand{\resp}{respectively\ }
\newcommand{\Subsection}{\subsection}
\providecommand{\nobreakdash}{\nobreak\hbox{-}}
\setlength{\emergencystretch}{2em}
\multlinegap0pt
\usepackage{amssymb}
\usepackage{enumerate}
\newtheorem{lemma}{Lemma}
\usepackage{cleveref}
\crefname{lemma}{Lemma}{Lemmas}
\title{The fractional $p$-Laplacian on hyperbolic spaces}
\author{Jongmyeong Kim, Minhyun Kim, Ki-Ahm Lee}
\date{Source: arXiv:2210.07029v3 (CC BY 4.0)}
\begin{document}
\maketitle

\noindent\emph{Source and licence.} The fractional $p$-Laplacian on hyperbolic spaces. Version arXiv:2210.07029v3 (CC BY 4.0), distributed by arXiv under the Creative Commons Attribution 4.0 International licence, http://creativecommons.org/licenses/by/4.0/, as recorded in the arXiv licence metadata for this exact version and shown on the abstract page https://arxiv.org/abs/2210.07029v3. Full licence notice: \texttt{licenses/arxiv-2210.07029-CC-BY-4.0.txt}.\par\smallskip

\noindent\textit{Editorial scope.} Counted unit: the article's lemma lem:zero (source0.tex lines 864--869). The article's complete original proof of it is delivered with it (source0.tex lines 871--1024, one \texttt{proof} environment of 154 lines). No mathematics is written, repaired or re-derived: the statement and the proof are the article's own lines, in the article's own notation, and no retained line is substituted or re-flowed.  Retained uncounted as the source-local context the proof needs: lines 248--253; lines 261--263; lines 265--267; lines 292--302; lines 322--328; lines 368--374.  Nothing was omitted from any retained span, so the package carries no omission record.  No retained line contains a citation, so the wrapper carries no bibliography and no bibliography entry.  No retained line contains a figure environment, an image-inclusion command or drawing code, so no image is delivered and the package declares no asset dependency.  Editorial formatting only, in addition: an AMS \texttt{amsart} preamble that restates the article's own declarations and macros used by the retained lines, namely the environments \texttt{lemma} on separate counters, together with the standard packages those lines need.  The retained environments print in this document as Lemma 1 in order of appearance, whereas the article prints the same items with the numbering of its own full document.  The complete native source of the article as distributed in the arXiv e-print for this exact version is bundled unchanged as \texttt{sources/132004/source0.tex}.
\begin{equation} \label{eq:K-asymptotic}
\begin{split}
&K_{\nu}(\rho) \sim \frac{1}{2} \Gamma(\nu) \left( \frac{\rho}{2} \right)^{-\nu} \quad\text{as } \rho \to 0^{+}, \text{ for } \nu > 0, \text{ and} \\
&K_{\nu}(\rho) \sim \sqrt{\frac{\pi}{2\rho}} e^{-\rho} \quad\text{as } \rho \to \infty.
\end{split}
\end{equation}

\begin{equation} \label{eq:K-scr}
\widetilde{K}_{\nu, a}(\rho) := \rho^{-\nu} K_{\nu}(a\rho).
\end{equation}

\begin{equation} \label{eq:scrK-recurrence}
-\partial_{\rho} (\widetilde{K}_{\nu, a}(f(\rho))) = a f'(\rho) f(\rho) \widetilde{K}_{\nu+1, a}(f(\rho))
\end{equation}

\begin{lemma} \label{lem:recurrence}
Let $\nu > 1/2$, $a \geq 1/2$, and $y \geq 0$. For $m \in \mathbb{N} \cup \lbrace 0 \rbrace$ define
\begin{equation*}
F_{m}(r) := \frac{\sinh r}{\sqrt{\cosh r - \cosh\rho}} \left( \frac{-\partial_{r}}{\sinh r} \right)^{m} \widetilde{K}_{\nu, a}\left( \sqrt{r^{2}+y^{2}} \right),
\end{equation*}
where $\widetilde{K}_{\nu, a}$ is the function given in \eqref{eq:K-scr}. Then, $F_{m}$ is integrable on $(\rho, \infty)$ and satisfies
\begin{equation} \label{eq:recurrence-F}
\left( \frac{-\partial_{\rho}}{\sinh \rho} \right) \int_{\rho}^{\infty} F_{m}(r) \,\mathrm{d}r = \int_{\rho}^{\infty} F_{m+1}(r) \,\mathrm{d}r
\end{equation}
for all $\rho>0$ and $m \in \mathbb{N} \cup \lbrace 0 \rbrace$.
\end{lemma}

\begin{lemma} \label{lem:positivity}
Let $\nu > 1/2$, $a \geq 1/2$, and $m \in \mathbb{N} \cup \lbrace 0 \rbrace$. Then, the function
\begin{equation*}
\rho \mapsto \left( \frac{-\partial_{\rho}}{\sinh \rho} \right)^{m} \widetilde{K}_{\nu,a}(\rho)
\end{equation*}
is positive.
\end{lemma}

\begin{lemma} \label{lem:even_dim}
  Let $a>0$ and $\nu > -\frac{n-1}{2}$. Then
  \begin{align*}
    \int^{\infty}_{\rho} \frac{\sinh^{-n/2+1} r}{\sqrt{\cosh r-\cosh \rho}} r^{-\nu}K_{n/2 + \nu} (ar) \,\mathrm{d} r \sim \sqrt{\frac{\pi}{2}} \frac{\Gamma(\nu+\frac{n-1}{2})}{\Gamma(\nu+\frac{n}{2})} \rho^{-\nu}\sinh^{-n/2 +1} (\rho)K_{n/2+\nu} (a\rho) 
  \end{align*}
  as $\rho \to 0^{+}$ up to dimensional constants.
\end{lemma}

\begin{lemma} \label{lem:zero}
Let $n\geq 2$ and $p > 1$. For any $R > 0$,
\begin{equation} \label{eq:zero}
\lim_{s \nearrow 1} c_{n, s, p} \int_{0}^{R} \rho^{p} \mathcal{K}_{n,s,p}(\rho)\sinh^{n-1}\rho \,\mathrm{d}\rho = \frac{1}{\pi^{\frac{n-1}{2}}} \frac{\Gamma(\frac{p+n}{2})}{\Gamma(\frac{p+1}{2})}.
\end{equation}
\end{lemma}

\begin{proof}
Let us first consider the odd-dimensional case $n=2m+1$ with $m \geq 1$. One can easily check that \eqref{eq:zero} is equivalent to
\begin{equation} \label{eq:odd}
\begin{split}
&\lim_{s \nearrow 1} (1-s) \int_{0}^{R} \rho^{p} \sinh^{2m} \rho \left( \frac{-\partial_{\rho}}{\sinh \rho} \right)^{m} \widetilde{K}_{\frac{1+sp}{2},m}(\rho) \, \mathrm{d}\rho \\
&= \frac{2^{m-1}}{p} \left( \frac{2}{m} \right)^{\frac{p+1}{2}} \Gamma\left( \frac{p+2m+1}{2} \right)
\end{split}
\end{equation}
by using $\lim_{s \nearrow 1} (1-s)|\Gamma(-s)| = 1$. Actually, we will prove the following statement, which is slightly stronger than \eqref{eq:odd}:
\begin{equation} \label{eq:zero-odd}
\begin{split}
\lim_{s \nearrow 1} (1-s) \int_{0}^{R} \rho^{p} \sinh^{2m} \rho \left( \frac{-\partial_{\rho}}{\sinh \rho} \right)^{m} \widetilde{K}_{\frac{1+sp}{2},a}(\rho) \, \mathrm{d}\rho \\
= \frac{2^{m-1}}{p} \left( \frac{2}{a} \right)^{\frac{p+1}{2}} \Gamma\left( \frac{p+2m+1}{2} \right) \quad\text{for each } a \geq 1.
\end{split}
\end{equation}
Let $\varepsilon \in (0,1)$, then there exists $\delta_{0} \in (0,1)$ such that
\begin{equation} \label{eq:delta-sinh}
1-\varepsilon \leq \frac{\sinh \rho}{\rho} \leq 1+\varepsilon
\end{equation}
for all $\rho \in (0, \delta_{0})$. Moreover, using the asymptotic behavior \eqref{eq:K-asymptotic} of the modified Bessel function, for each $s \in [0,1]$ we find $\delta_{s} > 0$ such that
\begin{equation} \label{eq:delta-K}
\frac{1-\varepsilon}{2} \Gamma\left( \frac{3+sp}{2} \right)\left( \frac{\rho}{2} \right)^{-\frac{3+sp}{2}} \leq K_{\frac{3+sp}{2}}(\rho) \leq \frac{1+\varepsilon}{2} \Gamma\left( \frac{3+sp}{2} \right)\left( \frac{\rho}{2} \right)^{-\frac{3+sp}{2}}
\end{equation}
for all $\rho \in (0, \delta_{s})$. Furthermore, since $\{K_{\nu}\}_{\nu \in [3/2, (3+p)/2]}$ is equicontinuous, we may assume that $\delta_{s}>0$ has been chosen continuously on $s$. Let us take $\delta = \delta_{0} \land \min_{s \in [0,1]} \delta_{s} \land R$, then $\delta = \delta(\varepsilon, p, R) > 0$, and \eqref{eq:delta-sinh} and \eqref{eq:delta-K} hold for all $\rho \in (0, \delta)$.

We fix $a \geq 1$ and denote by $G_{s, p, m, a}(\rho)$ the integrand in the left-hand side of \eqref{eq:zero-odd}. Then, $|G_{s, p, m, a}(\rho)|$ is bounded by the function $\sup_{0 \leq s \leq 1} |G_{s, p, m, a}(\rho)|$, which is independent of $s$ and bounded on a compact interval $[\delta/a, R]$. Thus, we have
\begin{equation*}
\lim_{s \nearrow 1} (1-s) \int_{\delta/a}^{R} G_{s, p, m, a}(\rho) \, \mathrm{d}\rho = 0,
\end{equation*}
and hence
\begin{equation*}
\lim_{s \nearrow 1} (1-s) \int_{0}^{R} G_{s, p, m, a}(\rho) \, \mathrm{d}\rho = \lim_{s \nearrow 1} (1-s) \int_{0}^{\delta/a} G_{s, p, m, a}(\rho) \, \mathrm{d}\rho.
\end{equation*}
Let us now prove \eqref{eq:zero-odd} by induction. When $m=1$, we first use \eqref{eq:scrK-recurrence} to have
\begin{equation*}
G_{s, p, 1, a}(\rho) = a\rho^{p-\frac{1+sp}{2}} K_{\frac{3+sp}{2}}(a\rho) \sinh \rho.
\end{equation*}
If $\rho < \delta/a$, then $\rho \leq a\rho < \delta$ since $a \geq 1$. Thus, we utilize \eqref{eq:delta-sinh} and \eqref{eq:delta-K} to obtain
\begin{equation*}
(1-\varepsilon)^2 \left( \frac{2}{a} \right)^{\frac{1+sp}{2}} \Gamma\left( \frac{3+sp}{2} \right) \rho^{p(1-s)-1} \leq G_{s, p, 1, a}(\rho) \leq (1+\varepsilon)^2 \left( \frac{2}{a} \right)^{\frac{1+sp}{2}} \Gamma\left( \frac{3+sp}{2} \right) \rho^{p(1-s)-1}.
\end{equation*}
This leads us to the inequalities
\begin{equation*}
\begin{split}
\lim_{s \nearrow 1} (1-s) \int_{0}^{R} G_{s, p, 1, a}(\rho) \, \mathrm{d}\rho
&= \lim_{s \nearrow 1} (1-s) \int_{0}^{\delta/a} G_{s, p, 1, a}(\rho) \, \mathrm{d}\rho \\
&\leq \lim_{s \nearrow 1} (1-s) (1+\varepsilon)^2 \left( \frac{2}{a} \right)^{\frac{1+sp}{2}} \Gamma\left( \frac{3+sp}{2} \right) \int_{0}^{\delta/a} \rho^{p(1-s)-1} \, \mathrm{d}\rho \\
&= (1+\varepsilon)^2 \frac{1}{p} \left( \frac{2}{a} \right)^{\frac{p+1}{2}} \Gamma\left( \frac{p+3}{2} \right)
\end{split}
\end{equation*}
and
\begin{equation*}
\lim_{s \nearrow 1} (1-s) \int_{0}^{R} G_{s, p, 1, a}(\rho) \, \mathrm{d}\rho \geq (1-\varepsilon)^{2} \frac{1}{p} \left( \frac{2}{a} \right)^{\frac{p+1}{2}} \Gamma\left( \frac{p+3}{2} \right).
\end{equation*}
Therefore, the statement \eqref{eq:zero-odd} for $m=1$ follows by taking $\varepsilon \to 0$.

Assume now that \eqref{eq:zero-odd} holds for $m \geq 1$. Then, a similar argument shows
\begin{equation*}
\begin{split}
&\lim_{s \nearrow 1} (1-s) \int_{0}^{R} G_{s, p, m+1, a}(\rho) \, \mathrm{d}\rho \\
&= \lim_{s \nearrow 1} (1-s) \int_{0}^{\delta/a} \rho^{p} \sinh^{2m+2} \rho \left( \frac{-\partial_{\rho}}{\sinh \rho} \right)^{m+1} \widetilde{K}_{\frac{1+sp}{2},a}(\rho) \, \mathrm{d}\rho \\
&\leq \lim_{s \nearrow 1} (1-s)(1+\varepsilon)^{2m+1} \int_{0}^{\delta/a} \rho^{p+2m+1}(-\partial_\rho) \left( \frac{-\partial_{\rho}}{\sinh \rho} \right)^{m} \widetilde{K}_{\frac{1+sp}{2},a}(\rho) \, \mathrm{d}\rho,
\end{split}
\end{equation*}
where nonnegativity of the integrands follows from \Cref{lem:positivity}. Using the integration by parts, \eqref{eq:delta-sinh}, and the induction hypothesis, we arrive at
\begin{equation*}
\begin{split}
&\lim_{s \nearrow 1} (1-s) \int_{0}^{R} G_{s, p, m+1, a}(\rho) \, \mathrm{d}\rho \\
&\leq (1+\varepsilon)^{2m+1} (p+2m+1) \lim_{s \nearrow 1} (1-s) \int_{0}^{\delta/a} \rho^{p+2m} \left( \frac{-\partial_\rho}{\sinh \rho} \right)^{m} \widetilde{K}_{\frac{1+sp}{2},a}(\rho) \, \mathrm{d}\rho \\
&\leq \frac{(1+\varepsilon)^{2m+1}}{(1-\varepsilon)^{2m}} (p+2m+1) \lim_{s \nearrow 1} (1-s) \int_{0}^{\delta/a} \rho^{p} \sinh^{2m}\rho \left( \frac{-\partial_\rho}{\sinh \rho} \right)^{m} \widetilde{K}_{\frac{1+sp}{2},a}(\rho) \, \mathrm{d}\rho \\
&= \frac{(1+\varepsilon)^{2m+1}}{(1-\varepsilon)^{2m}} (p+2m+1) \frac{2^{m-1}}{p} \left( \frac{2}{a} \right)^{\frac{p+1}{2}} \Gamma\left( \frac{p+2m+1}{2} \right) \\
&= \frac{(1+\varepsilon)^{2m+1}}{(1-\varepsilon)^{2m}} \frac{2^{m}}{p} \left( \frac{2}{a} \right)^{\frac{p+1}{2}} \Gamma\left( \frac{p+2m+3}{2} \right).
\end{split}
\end{equation*}
Similarly, we obtain
\begin{equation*}
\lim_{s \nearrow 1} (1-s) \int_{0}^{R} G_{s, p, m+1, a}(\rho) \, \mathrm{d}\rho \geq \frac{(1-\varepsilon)^{2m+1}}{(1+\varepsilon)^{2m}} \frac{2^{m}}{p} \left( \frac{2}{a} \right)^{\frac{p+1}{2}} \Gamma\left( \frac{p+2m+3}{2} \right),
\end{equation*}
from which \eqref{eq:zero-odd} for $m+1$ follows by taking $\varepsilon \to 0$. The statement \eqref{eq:zero-odd} has been proved for all $m \in \mathbb{N}$, finishing the proof of \eqref{eq:zero} for the odd-dimensional case.

Let us next consider the even-dimensional case $n = 2m$ with $m \geq 1$. In this case, \eqref{eq:zero} is equivalent to
\begin{equation*}
\begin{split}
&\lim_{s \nearrow 1} (1-s) \int_{0}^{R} \rho^{p} \sinh^{2m-1}\rho \int_{\rho}^{\infty} \frac{\sinh r}{\sqrt{\cosh r - \cosh \rho}} \left( \frac{-\partial_{r}}{\sinh r} \right)^{m} \widetilde{K}_{\frac{1+sp}{2},\frac{2m-1}{2}}(r) \, \mathrm{d}r \,\mathrm{d}\rho \\
&= \sqrt{\frac{\pi}{2}} \frac{2^{m-1}}{p} \left( \frac{2}{m-1/2} \right)^{\frac{p+1}{2}} \Gamma\left(\frac{p+2m}{2} \right).
\end{split}
\end{equation*}
As in the odd-dimensional case, we will prove a stronger statement:
\begin{equation} \label{eq:zero-even}
\begin{split}
&\lim_{s \nearrow 1} (1-s) \int_{0}^{R} \rho^{p} \sinh^{2m-1}\rho \int_{\rho}^{\infty} \frac{\sinh r}{\sqrt{\cosh r - \cosh \rho}} \left( \frac{-\partial_{r}}{\sinh r} \right)^{m} \widetilde{K}_{\frac{1+sp}{2},a}(r) \, \mathrm{d}r \,\mathrm{d}\rho \\
&= \sqrt{\frac{\pi}{2}} \frac{2^{m-1}}{p} \left( \frac{2}{a} \right)^{\frac{p+1}{2}} \Gamma\left(\frac{p+2m}{2} \right) \quad\text{for each } a \geq 1/2.
\end{split}
\end{equation}
Recall that we have taken $\delta$ so that \eqref{eq:delta-sinh} and \eqref{eq:delta-K} hold for all $\rho \in (0, \delta)$. Let us fix $a \geq 1/2$. By \Cref{lem:even_dim}, for each $s \in [0,1]$ we find $\tilde{\delta}_{s} > 0$ such that
\begin{equation} \label{eq:delta-intK}
\begin{split}
(1-\varepsilon) \sqrt{\frac{\pi}{2}} \frac{\Gamma(\frac{2+sp}{2})}{\Gamma(\frac{3+sp}{2})} \rho^{-\frac{1+sp}{2}} K_{\frac{3+sp}{2}}(a\rho)
&\leq \int_{\rho}^{\infty} \frac{r^{-\frac{1+sp}{2}} K_{\frac{3+sp}{2}}(ar)}{\sqrt{\cosh r - \cosh \rho}} \, \mathrm{d}r \\
&\leq (1+\varepsilon) \sqrt{\frac{\pi}{2}} \frac{\Gamma(\frac{2+sp}{2})}{\Gamma(\frac{3+sp}{2})} \rho^{-\frac{1+sp}{2}} K_{\frac{3+sp}{2}}(a\rho)
\end{split}
\end{equation}
for all $\rho \in (0, \tilde{\delta}_{s})$. Moreover, we may assume that $\tilde{\delta}_{s}$ has been chosen continuously on $s$ since $\{K_\nu\}_{\nu\in [3/2,(3+p)/2]}$ is equicontinuous. Let $\tilde{\delta} = \delta \land \min_{s \in [0,1]} \tilde{\delta}_{s}$, then $\delta = \delta(\varepsilon, p, R, a) > 0$ and \eqref{eq:delta-intK} holds for all $\rho \in (0, \tilde{\delta})$.

We denote by $H_{s, p, m, a}(\rho)$ the integrand in the left-hand side of \eqref{eq:zero-even}. Then, the same argument as in the odd-dimensional case shows
\begin{equation*}
\lim_{s \nearrow 1} (1-s) \int_{0}^{R} H_{s, p, m, a}(\rho) \, \mathrm{d}\rho = \lim_{s \nearrow 1} (1-s) \int_{0}^{\frac{\delta}{2a}} H_{s, p, m, a}(\rho) \, \mathrm{d}\rho.
\end{equation*}
We argue by induction again to prove \eqref{eq:zero-even}. If $m=1$, then
\begin{equation*}
H_{s, p, 1, a}(\rho) = a \rho^{p} \sinh\rho \int_{\rho}^{\infty} \frac{1}{\sqrt{\cosh r - \cosh \rho}} r^{-\frac{1+sp}{2}} K_{\frac{3+sp}{2}}(ar) \, \mathrm{d}r.
\end{equation*}
Since $a \geq 1/2$, we have $\rho < \delta$ and $a\rho < \delta$ for $\rho < \frac{\delta}{2a}$. Thus, by \eqref{eq:delta-sinh}, \eqref{eq:delta-intK}, and \eqref{eq:delta-K}, we obtain
\begin{equation*}
\begin{split}
(1-\varepsilon)^{3} \sqrt{\frac{\pi}{2}} \left( \frac{2}{a} \right)^{\frac{1+sp}{2}} \Gamma\left( \frac{2+sp}{2} \right) \rho^{p(1-s)-1}
&\leq H_{s, p, 1, a}(\rho) \\
&\leq (1+\varepsilon)^{3} \sqrt{\frac{\pi}{2}} \left( \frac{2}{a} \right)^{\frac{1+sp}{2}} \Gamma\left( \frac{2+sp}{2} \right) \rho^{p(1-s)-1}.
\end{split}
\end{equation*}
Therefore, we have
\begin{equation*}
\begin{split}
(1-\varepsilon)^{3} \sqrt{\frac{\pi}{2}} \frac{1}{p} \left( \frac{2}{a} \right)^{\frac{p+1}{2}} \Gamma\left( \frac{p+3}{2} \right)
&\leq \lim_{s \nearrow 1} (1-s) \int_{0}^{R} H_{s, p, 1, a}(\rho) \, \mathrm{d}\rho \\
&\leq (1+\varepsilon)^{3} \sqrt{\frac{\pi}{2}} \frac{1}{p} \left( \frac{2}{a} \right)^{\frac{p+1}{2}} \Gamma\left( \frac{p+3}{2} \right),
\end{split}
\end{equation*}
from which we deduce \eqref{eq:zero-even} for $m=1$ by taking $\varepsilon \to 0$.

Suppose that \eqref{eq:zero-even} is true for $m \geq 1$. Then, by \eqref{eq:delta-sinh}, \Cref{lem:recurrence}, and \Cref{lem:positivity}, we have
\begin{equation*}
\begin{split}
&\lim_{s \nearrow 1} (1-s) \int_{0}^{R} H_{s, p, m+1, a}(\rho) \, \mathrm{d}\rho \\
&= \lim_{s \nearrow 1} (1-s) \int_{0}^{\frac{\delta}{2a}} \rho^{p} \sinh^{2m+1}\rho \int_{\rho}^{\infty} \frac{\sinh r}{\sqrt{\cosh r - \cosh \rho}} \left( \frac{-\partial_{r}}{\sinh r} \right)^{m+1} \widetilde{K}_{\frac{1+sp}{2},a}(r) \, \mathrm{d}r \, \mathrm{d}\rho \\
&\leq \lim_{s \nearrow 1} (1-s) (1+\varepsilon)^{2m} \int_{0}^{\frac{\delta}{2a}} \rho^{p+2m} (-\partial_{\rho}) \int_{\rho}^{\infty} \frac{\sinh r}{\sqrt{\cosh r - \cosh \rho}} \left( \frac{-\partial_{r}}{\sinh r} \right)^{m} \widetilde{K}_{\frac{1+sp}{2},a}(r) \, \mathrm{d}r \, \mathrm{d}\rho.
\end{split}
\end{equation*}
Using the integration by parts, \eqref{eq:delta-sinh}, and the induction hypothesis, we arrive at
\begin{equation*}
\begin{split}
&\lim_{s \nearrow 1} (1-s) \int_{0}^{R} H_{s, p, m+1, a}(\rho) \, \mathrm{d}\rho \\
&= (1+\varepsilon)^{2m} (p+2m) \\
&\quad\times \lim_{s \nearrow 1} (1-s) \int_{0}^{\frac{\delta}{2a}} \rho^{p+2m-1} \int_{\rho}^{\infty} \frac{\sinh r}{\sqrt{\cosh r - \cosh \rho}} \left( \frac{-\partial_{r}}{\sinh r} \right)^{m} \widetilde{K}_{\frac{1+sp}{2},a}(r) \, \mathrm{d}r \, \mathrm{d}\rho \\
&\leq \frac{(1+\varepsilon)^{2m}}{(1-\varepsilon)^{2m-1}} (p+2m) \lim_{s \nearrow 1} (1-s) \int_{0}^{\frac{\delta}{2a}} H_{s, p, m, a}(\rho) \, \mathrm{d}\rho \\
&= \frac{(1+\varepsilon)^{2m}}{(1-\varepsilon)^{2m-1}} \sqrt{\frac{\pi}{2}} \frac{2^{m}}{p} \left( \frac{2}{a} \right)^{\frac{p+1}{2}} \Gamma\left(\frac{p+2m+2}{2} \right).
\end{split}
\end{equation*}
The inequality
\begin{equation*}
\lim_{s \nearrow 1} (1-s) \int_{0}^{R} H_{s, p, m+1, a}(\rho) \, \mathrm{d}\rho \geq \frac{(1-\varepsilon)^{2m}}{(1+\varepsilon)^{2m-1}} \sqrt{\frac{\pi}{2}} \frac{2^{m}}{p} \left( \frac{2}{a} \right)^{\frac{p+1}{2}} \Gamma\left( \frac{p+2m+2}{2} \right)
\end{equation*}
can be obtained in the same way. Thus, we conclude that \eqref{eq:zero-even} for $m+1$ holds by taking $\varepsilon \to 0$. This finishes the proof for the even-dimensional case.
\end{proof}

\end{document}
